%% ============================================================================
%% SR5 Level-Constrained Planning -- conference paper (IEEE conference format)
%%
%% Target venue: IEEE ICRA / IROS (IEEEtran conference style, US-letter, 2-column).
%% The IEEEtran class + BibTeX style live under ./IEEEtran/ (official CTAN package).
%%
%% Build:  see README.md in this directory (latexmk -pdf main.tex).
%% ============================================================================
\documentclass[conference]{IEEEtran/IEEEtran}
\IEEEoverridecommandlockouts
% The preceding line is only needed to identify funding in the first footnote.

\usepackage{cite}
\usepackage{amsmath,amssymb,amsfonts}
\usepackage{algorithmic}
\usepackage{graphicx}
\usepackage{textcomp}
\usepackage{xcolor}
\usepackage{booktabs}
\usepackage{subcaption}

% hyperref should be loaded last.
\usepackage[hidelinks]{hyperref}

\graphicspath{{figures/}}

% Convenience macros for constrained-planning notation.
\newcommand{\qvec}{\mathbf{q}}
\newcommand{\Tgoal}{T_{\text{goal}}}
\newcommand{\zee}{z_{\text{ee}}}
\newcommand{\zw}{z_{w}}

\begin{document}

\title{Level-Constrained Trajectory Planning for Manipulators\\
via Seed Construction, CuRobo Repair, and a Closed-Loop\\
Learning--Optimization System\\
\thanks{This work was carried out on the SR5 manipulator platform.
Corresponding author: (email).}}

\author{\IEEEauthorblockN{First Author}
\IEEEauthorblockA{\textit{Affiliation} \\
\textit{Institution}\\
City, Country \\
email@domain}
\and
\IEEEauthorblockN{Second Author}
\IEEEauthorblockA{\textit{Affiliation} \\
\textit{Institution}\\
City, Country \\
email@domain}
}

\maketitle

\begin{abstract}
Keeping a manipulator's end effector level throughout an entire path---the
``carrying a tray of water'' constraint---turns point-to-point planning into
motion on a low-dimensional constraint manifold, where feasible solutions are
scarce and gradient-based optimizers are highly sensitive to their initial
seed. We present a level-constrained trajectory planning system for the SR5
manipulator that (i) constructs a diverse family of rule-based level-preserving
seeds via horizontal-pose interpolation and branch-consistent inverse
kinematics, (ii) repairs and refines every seed with GPU-batched CuRobo
trajectory optimization, and (iii) selects the executable trajectory with a
level-first hard validator. We further wrap this planner in a closed-loop
learning--optimization system---\emph{data generation, model learning,
optimization validation, failure fallback, data update}---in which a diffusion
seed model and a success critic learn the distribution of CuRobo-repairable
seeds, while every learned candidate still returns to CuRobo repair, hard
validation, and rule-based fallback. We report the constrained-planning success
rate under a fixed time budget and analyze the contribution of each component.
\end{abstract}

\begin{IEEEkeywords}
constrained motion planning, manipulation, trajectory optimization, diffusion
models, inverse kinematics, learning for planning
\end{IEEEkeywords}

\section{Introduction}
\label{sec:intro}
% Motivation: end-effector level (orientation-path) constraint; "carry water" tasks
% (trays, welding, spraying, polishing). Distinguish from point-to-point planning.
% Contribution list:
%  1) A multi-family rule-based level seed constructor with branch-consistent IK.
%  2) A CuRobo repair + level-first hard-validation pipeline.
%  3) A closed-loop learning--optimization system that treats every candidate
%     (success, failed seed, fallback recovery) as first-class training data.
Placeholder.

\section{Related Work}
\label{sec:related}
% Sampling-based constrained planning (TSR/CBiRRT, OMPL constrained, cpRRTC, McVAMP).
% IK-based methods (IKFast, TRAC-IK, RelaxedIK, IKFlow).
% Trajectory optimization (CHOMP, STOMP, TrajOpt, GPMP, CuRobo).
% Learning-based seeding (DiffusionSeeder, Motion Planning Diffusion, PRESTO).
Placeholder.

\section{Problem Formulation}
\label{sec:problem}
% Trajectory as (N+1) x n joint-configuration sequence P = [q_0, ..., q_N].
% Objective: smoothness (velocity/acceleration/jerk) + regularization.
% Constraints: start q_0 = q_start; goal T_ee(q_N) = T_goal;
%   level path constraint z_ee(q_k)^T z_w >= cos(eps) for all k;
%   joint limits; velocity/acceleration; collision d(q_k, O) >= d_safe;
%   joint-space continuity ||q_{k+1} - q_k|| <= dq_max.
The level constraint at each waypoint is expressed as
\begin{equation}
  \zee(\qvec_k)^\top \zw \ge \cos\epsilon, \qquad k = 0,1,\dots,N,
  \label{eq:level}
\end{equation}
where $\zee(\qvec_k)$ is the tool axis and $\zw$ the target world direction.

\section{Method}
\label{sec:method}

\subsection{Multi-Family Level Seed Construction}
\label{sec:seeds}
% Cartesian linear interpolation of position; pose restricted to the horizontal
% family via goal-relative twist interpolation Q_i = Q_g * T_y(theta_i);
% branch-consistent IK selection minimizing joint jump + trend + goal-anchor +
% limit-margin costs; smoothed / bridged variants; segmented optimization.

\subsection{CuRobo Repair and Hard Validation}
\label{sec:repair}
% Seeds enter trajopt_solver.solve_pose(seed_traj=...); level-first candidate
% selection ranks by alignment validity, start-gap, joint-step jump, twist
% smoothness, path cost; hard validator checks alignment/goal/continuity/
% joint-limit/velocity-acceleration.

\subsection{Closed-Loop Learning--Optimization System}
\label{sec:closedloop}
% Diffusion seed model learns p(tau_seed | q_s, T_g, O, C); success critic
% predicts (P_success, alignment risk, collision risk, opt cost); every learned
% candidate returns to CuRobo repair + hard validation; failures fall back to
% rule seeds; all outcomes feed the next dataset version.

\section{Experiments}
\label{sec:experiments}
% Main metric: success rate after CuRobo repair + hard validation under a fixed
% total time budget. Ablations: rule-only vs diffusion-only vs diffusion+critic
% vs mixed fallback. Report SR5 real-robot and simulation stress tests.

\begin{table}[t]
\caption{Constrained-planning success rate (placeholder).}
\label{tab:main}
\centering
\begin{tabular}{lcc}
\toprule
Method & Success rate & Mean solve time (s) \\
\midrule
CuRobo native (no seed) & -- & -- \\
Rule seeds + CuRobo + selection & -- & -- \\
Diffusion seeds only & -- & -- \\
Diffusion + critic & -- & -- \\
Mixed fallback & -- & -- \\
\bottomrule
\end{tabular}
\end{table}

\section{Conclusion}
\label{sec:conclusion}
Placeholder.

\section*{Acknowledgment}
Placeholder.

\bibliographystyle{IEEEtran/bibtex/IEEEtran}
\bibliography{references}

\end{document}
